%include all the necessaru packages required to compile the main document (main.tex) 

\usepackage[portuguese]{babel}  %Define portuguese as the document language 
\usepackage{graphicx}   %Required for inserting images
\usepackage{amsmath}    %Required for math environment 
\numberwithin{equation}{section}
\usepackage{xcolor}%    %Required for colored text 
\usepackage{amssymb}
\usepackage[utf8]{inputenc}
\usepackage{hyperref}
\hypersetup{
	colorlinks=true,
	linkcolor=blue,
	filecolor=magenta,      
	urlcolor=cyan,
	pdftitle={Overleaf Example},
	pdfpagemode=FullScreen,
}
\usepackage{booktabs} % For better table rules
\usepackage{caption,subcaption}
\usepackage{amsfonts}
\usepackage{algorithmic}
\usepackage{amsthm}
\usepackage{tikz}
\usepackage{pgfplots}
\pgfplotsset{compat=1.16}

%\usepackage[a4paper]{geometry}
%%%%%%%%%%%%%%%%%%%%%%%%%%%%%%%%%%%%%%%%%%%%%%%%%%%%%%%%%%%%%%%%%%%%%%%%%%

\newcommand{\propositionautorefname}{Proposição}% naming the autoref proposition
\newcommand{\exempleautorefname}{Exemplo}       % naming the autoref exemple
\newcommand{\corolaryautorefname}{Corolário}    % naming the autoref corolary
\newcommand{\definitionautorefname}{Definição}  % naming the autoref definition
\newcommand{\lemmaautorefname}{Lema}            % naming the autoref lemma
\newcommand{\exerciseautorefname}{Exercício}    % naming the autoref lemma
\renewcommand{\theoremautorefname}{Teorema}     % renaming the autoref theorem

\renewcommand{\appendixname}{Apêndice}          % renaming the body "appendix"
\renewcommand{\appendixautorefname}{Apêndice}   % renaming the autoref appendix
\renewcommand*{\proofname}{Demonstração}        % renaming the body "proof"

\setlength{\parskip}{0.2cm}                     % control the space between paragrhaphes

\setlength{\skip\footins}{1.5\baselineskip}      % control the space over the line of footnote

\newtheoremstyle{style-one}          % name of the style to be used
{10pt}                             % measure of space to leave above the theorem
{8pt}                              % measure of space to leave below the theorem 
{\slshape}                         % name of font to use in the body of the theorem
{15pt}                             % measure of space to indent
{\bfseries}                        % name of head font
{.}                                % punctuation between head and body
{ }                                % space after theorem head
{\thmname{#1}\thmnumber{ #2}{\thmnote{ (#3)}}}

\theoremstyle{style-one}
\newtheorem{definition}{Definição}[section]
\newtheorem{proposition}{Proposição}[section]
\newtheorem{theorem}{Teorema}[section]
\newtheorem{exemple}{Exemplo}[section]
\newtheorem{exercise}{Exercício}[section]
\newtheorem{lemma}{Lema}[section]

\newtheorem*{definition*}{Definição}
\newtheorem*{corolary}{Corolário}

\renewcommand\qedsymbol{$\blacksquare$}



%%%%%%%%%%%%%%%%%%%%%%%%%%%%%%%%%%%%%%%%%%%%%%%%%%%%%%%%%%%%%%%%%%%%%%%%%%
